% Generated draft: requires mathematical and source review.
\documentclass[11pt]{article}
\usepackage[margin=25mm]{geometry}
\usepackage{fontspec}
\setmainfont{Latin Modern Roman}
\usepackage{amsmath,amssymb,mathtools,unicode-math}
\setmathfont{Latin Modern Math}
\usepackage{enumitem,needspace,xurl,hyperref}
\setlength{\parindent}{0pt}
\setlength{\parskip}{5pt}
\setlength{\emergencystretch}{3em}
\newcommand{\N}{\mathbb{N}}
\newcommand{\Z}{\mathbb{Z}}
\newcommand{\Q}{\mathbb{Q}}
\newcommand{\R}{\mathbb{R}}
\newcommand{\C}{\mathbb{C}}
\newcommand{\F}{\mathbb{F}}
\newcommand{\A}{\mathbb{A}}
\newcommand{\PP}{\mathbb{P}}
\newcommand{\E}{\mathbb{E}}
\newcommand{\eps}{\varepsilon}
\newcommand{\dd}{\,\mathrm{d}}
\begin{document}
\section*{UnsolvedMath ID: 20001980}
\textbf{Source problem code:} catalogue entry
\subsection{Definitions and mathematical statement}
M. Libine's Question \(11.2\) asks how to parametrize symplectic forms on \(T^*M\) ``up to this equivalence'' if ``such a diffeomorphism'' does not exist. Those phrases refer to the preceding Question \(11.1\), which is not reproduced. The aid reconstructs the equivalence as a diffeomorphism preserving the zero section setwise and identifying the forms up to sign. This is stated explicitly, rather than assumed to be the only possible cotangent-bundle equivalence.

For the aid's fully specified special case, write
\[
 T^*S^1=(\R/2\pi\Z)_q\times\R_p,\qquad Z=\{p=0\}.
\]
For every smooth \(f:\R\to(0,\infty)\), consider
\[
 \omega_f=f(p)\,\dd p\wedge\dd q
      =\dd(F_f(p)\,\dd q),\qquad F_f(p)=\int_0^p f(s)\,\dd s.
\]
The indicated primitive pulls back to zero on \(Z\). Say \(\omega_f\sim\omega_g\) if there exists a diffeomorphism \(\Phi:T^*S^1\to T^*S^1\) with \(\Phi(Z)=Z\) and \(\Phi^*\omega_g=\omega_f\) or \(-\omega_f\). The diffeomorphism need not preserve individual cotangent fibers or fix \(Z\) pointwise.

Define the positive extended end areas
\[
 A_-(f)=2\pi\int_{-\infty}^0f(p)\,\dd p,\qquad
 A_+(f)=2\pi\int_0^\infty f(p)\,\dd p
       \quad\in(0,\infty].
\]
The aid's classification theorem is
\[
 \omega_f\sim\omega_g
 \quad\Longleftrightarrow\quad
 \{A_-(f),A_+(f)\}=\{A_-(g),A_+(g)\},
\]
where pairs are unordered with multiplicity. Every such pair occurs, so the resulting family quotient is
\[
 ((0,\infty]\times(0,\infty])/\mathfrak S_2.
\]
Finite and infinite ends are both included. For equal ordered end areas, \(F_f,F_g\) have the same image interval and the fiber-coordinate change \(p\mapsto F_g^{-1}(F_f(p))\) gives an explicit equivalence. Reversing ends handles the unordered case under the signed convention.

The canonical form \(\dd p\wedge\dd q\) has pair \(\{\infty,\infty\}\), so a finite end obstructs equivalence to it in this category. Total area alone is incomplete relative to the marked zero section: finite pairs \(\{1,3\}\) and \(\{2,2\}\), for example, have the same total and distinct invariants. This classification concerns fiber-dependent forms on \(T^*S^1\), not all symplectic forms there or on arbitrary cotangent bundles.

Comment \(11.3\), also Libine, motivates the general question by extensions of the Berline--Vergne localization formula to noncompact reductive group actions; no such localization extension is proved by the end-area calculation. The next heading, ``\(12\) Volume growth of hyperk\"ahler manifolds,'' is spillover without a further question in this record.
\subsection{Short English statement}
The general cotangent-bundle equivalence depends on missing preceding context. In the aid's precise cylinder family, the complete invariant is the unordered pair of areas on the two sides of the preserved zero section, allowing infinite areas. This yields a special moduli slice and obstructions to the canonical form, rather than a general classification.
\subsection{Sources}
\begin{itemize}
\item Local UnsolvedMath registry, permanent ID \(20001980\), source code \texttt{catalogue entry}; canonical collection \texttt{aim-geometry-notes.json}, notes[317].
\item AIM workshop ``Moment maps and surjectivity in various geometries,'' Question \(11.2\) and Comment \(11.3\), both M. Libine: \url{https://aimath.org/WWN/momentmaps/momentmaps.pdf}. The equivalence refers to missing Question \(11.1\); Section \(12\)'s heading is appended spillover. No PDF was opened in this drafting session.
\item The comment names the Berline--Vergne localization formula, extracted as ``Berligne-Vergne,'' and extensions to noncompact reductive actions, without a full bibliographic entry. Original remarks and literature are empty; proposer/year registry fields are null despite Libine's attribution.
\item Supplied research attempt \(1\), \texttt{partial\_result}, run status \texttt{partially\_solved}, rigor label ``proved,'' novelty confidence ``low.'' Registry classification \texttt{PARTIAL-PROGRESS}, updated 14 August 2026.
\end{itemize}
\end{document}
